\scanpage{163}
\mspage{380}

here these \unsure{delightful} \unsure{scales} of linear algebra,
or even of scheme-theoretic interpretations of the kind that the scheme of solutions of (*)
(for $f,\alpha_0,\mu$ variable) is a smooth \inserted{relative} scheme of relative dimension 4, \unsure{which etc.}

I will only spell out, in the present adelic context, manifestly \unsure{too} special, the

\uline{Proposition} Let $\deleted{\unsure{a,b}}c,d\quad\deleted{\unsure{r,s}}t,u\in\widehat{\mathbb Z}$,
for them to \unsure{correspond} to data
$\alpha_0=\begin{pmatrix}a&b\\c&d\end{pmatrix}$,
$f=\begin{pmatrix}r&s\\t&u\end{pmatrix}$ satisfying
\[
\underbrace{ad-bc}_{\substack{\unsure{\text{definition}}\\=\varepsilon_3}}\in\{\pm1\},
\quad ru-st=1,
\quad\alpha_0\rho(\alpha_0)^{-1}=\varepsilon_2\ell_0^{\nu}(f\sigma_\infty(\unsure{f})^{-1})
\]
with $\varepsilon_2\in\{\pm1\}$ and $\unsure{\nu}\in\widehat{\mathbb Z}$
\unsure{suitable|suitable} (cf. (48) (49) (53)),
it \unsure{f|is necessary} and \deleted{\illegible{1}} it suffices that one have
\[
\tag{85}
c^2+\unsure{cd}+d^2=\pm(t^2+u^2)\quad\in\widehat{\mathbb Z}^*
\]
Then \deleted{\unsure{signs} $\varepsilon_2,\varepsilon_3$ and the scalar $\unsure{\nu}\in\widehat{\mathbb Z}$ are}
\inserted{$(\varepsilon_2,\varepsilon_3,2\nu+1)\in\{\pm1\}\times\{\pm1\}\times\widehat{\mathbb Z}^*$ is determined}
modulo sign (i.e. modulo passing to: $(-\varepsilon_2,-\varepsilon_3,-\nu)$)
\deleted{then uniquely determined as follows: One}
\inserted{by $a,b,c,d$, by the \deleted{condition} relation}
\inserted{(where $\mu=2\nu+1$)}
\deleted{defines \illegible{1}}
\[
\deleted{\varepsilon_3\unsure{\equiv|=}\unsure{(c^2+cd+d^2)}\unsure{\in\widehat{\mathbb Z}^*}}
\]
\[
\tag{86}
\frac1\mu=\varepsilon_2(t^2+u^2)=\varepsilon_3(c^2+cd+u^2)
\]
\margin{\unsure{by}}
\[
\deleted{\unsure{\mu|\frac1\mu}=\unsure{\varepsilon_3}(c^2+cd+d^2)=}
\]
\deleted{\illegible{3}}
\inserted{(which determines $\varepsilon_3,\mu$ when one has chosen $\varepsilon_2$)}
and one has
\[
\tag{87}
\left\{\begin{aligned}
\varepsilon_2&\equiv\mu\quad(4)\\
\varepsilon_3&\equiv\mu\quad(3)\\
\nu&=(\mu-1)/2
\end{aligned}\right.
\]

\margin{162}

\scanpage{164}

The solutions are then obtained as follows: one chooses
$\varepsilon_2,\varepsilon_3,\unsure{\mu|\mu_2}$ as above (two possibilities), \deleted{then} and arbitrarily $r,s\in\widehat{\mathbb Z}$ satisfying
\[
ur-ts=1
\]
(there are always some (for $r^2+s^2\in\widehat{\mathbb Z}^*$ --- hence
$\forall p\in\mathbb P$, $r_p$ and $s_p$ are not both in $p\mathbb Z_p$)
and if $r_0,s_0$ is a solution, the others are obtained in parametric form in $\lambda$
\[
\tag{87}
(r,s)=(r_0,s_0)\unsure{+|-}\lambda(\unsure{t|u},\unsure{u|t})
=(r_0\unsure{+|-}\lambda t,s_0\unsure{+|-}\lambda u).
\]
Having chosen $r,s$ hence $f=\begin{pmatrix}r&s\\t&u\end{pmatrix}\in\mathrm{Sl}(2,\widehat{\mathbb Z})$,

\deleted{
one deduces from this
\[
f\sigma_\infty(f)^{-1}=f\,{}^tf
=\begin{pmatrix}R&S\\T&U\end{pmatrix}
\]
where $R=\unsure{r}^2+s^2$, $S=T=rt+su$, $U=t^2+u^2$, and one successively determines the quantities
\[
\tag{88}
\left\{\begin{aligned}
D&=1/\mu=\varepsilon_3(c^2+cd+d^2)
=\deleted{\varepsilon_2\underbracket{\unsure{U}}_{\inserted{t^2+}}=\varepsilon_2(\unsure{t^2}+\illegible{1})}
\inserted{\varepsilon_2(t^2+u^2)=\varepsilon_2U}\\
C&=\varepsilon_2T\ (=\varepsilon_2S)\\
B&=C+D-1=\varepsilon_2(S+U)-1\\
A&=\frac1D(1+BC)=\frac1{\varepsilon_2U}(1+S(S+U)-\varepsilon_2S)
\end{aligned}\right.
\]
\margin{$\unsure{t^2+}$}
whence a matrix $\begin{pmatrix}A&B\\C&D\end{pmatrix}$ satisfying
$C+D-B=1$ and $D\in\widehat{\mathbb Z}^*$, and satisfying relation (57)
\[
\begin{pmatrix}A+2\nu C&B+2\nu D\\C&D\end{pmatrix}
=\varepsilon_2\begin{pmatrix}R&S\\T&U\end{pmatrix}
\]
which here reduces \unsure{to} the relation in the upper left corner
$A+2\nu C=\varepsilon_2R$ i.e.
}

\margin{163}

\scanpage{165}
\mspage{381}

it remains to choose $a,b$ so as to satisfy the two relations
$\det\begin{pmatrix}a&b\\c&d\end{pmatrix}=\varepsilon_3$
and $C=\varepsilon_2 T$ (\unsure{cf.} (61) (62))
\inserted{i.e.}
\[
\tag{88}
\left\{\begin{aligned}
d\uline{a}-c\uline{b}&=\varepsilon_3\\
(\unsure{c}+d)\uline{a}+d\uline{b}
&=\varepsilon_2\varepsilon_3
\underbrace{(\uline{r}t+\uline{s}u)}_{S}
+\underbrace{(c^2+cd+d^2)}_{1/\mu\varepsilon_3}
\end{aligned}\right.
\]
which is a system of two linear equations in $\uline{a},\uline{b}$,
whose determinant is as if by chance
$\unsure{c^2}+cd+d^2=1/\mu\varepsilon_3\in\widehat{\mathbb Z}^{*}$
--- which ensures a unique solution!

It seems more natural perhaps to start from the multiplier $\mu$, and to
look for the systems
$\begin{pmatrix}a&b\\c&d\end{pmatrix}$
$\begin{pmatrix}r&s\\t&u\end{pmatrix}$ corresponding to it,
but putting the emphasis on $(c,d)\ (t,u)$ in the first place.
Using $\mu$ one determines $\varepsilon_2,\varepsilon_3$ by
\[
\mu\equiv\varepsilon_2\ (2),\qquad
\mu\equiv\varepsilon_3\ (3),
\]
\deleted{\unsure{then}} which allows one to \deleted{\unsure{define}} represent
$1/\mu$ in the form
\[
1/\mu=\varepsilon_2(t^2+u^2)=\varepsilon_3(c^2+cd+d^2)
\]
(with $t,u,c,d\in\widehat{\mathbb Z}$, the systems $(t,u)$ as well as
$(c,d)$ ``\unsure{depend} on $1$ parameter'', \inserted{each}
in a very concrete sense that one could probably make precise if one cared to,
\deleted{\unsure{a wide choice}}\inserted{with}
explicit parameters \deleted{\illegible{1}} $\sigma,\tau$\ldots{} Once
$c,d,\unsure{t,u}$ have been chosen, one chooses $r,\unsure{s|t}$ as has

\margin{164}

\scanpage{166}

been said (with a parameter $\unsure{\lambda}$ \inserted{entering}
linearly), which then determines $r,s$ uniquely, in a way
that involves $\lambda$ \deleted{linearly} in a
linear-affine way.

In fact, for $c,d,t,u$ \inserted{and $\varepsilon_2,\varepsilon_3$}
fixed, \deleted{the equations}\inserted{so as to}
satisfy: $\varepsilon_2(t^2+u^2)=\varepsilon_3(c^2+cd+d^2)$,
the three equations relating $a,b,r,s$ are \unsure{linear-affine}
\[
\tag{89}
\left\{\begin{aligned}
d\uline{a}-c\uline{b}&=\varepsilon_3\\
u\uline{r}-t\uline{s}&=1\\
(c+d)\uline{a}+d\uline{b}
-\varepsilon_2\varepsilon_3(u\uline{s}+t\uline{r})&=c^2+cd+d^2
\end{aligned}\right.
\]
\margin{we have underlined the ``unknown coefficients''
$\uline{a},\uline{b},\uline{r},\uline{s}$ to be determined}
and the \deleted{determinant of this} system is ``of rank $3$''
\deleted{\unsure{det}}
--- the set of its solutions is an ``affine space of dimension $1$''
over $\widehat{\mathbb Z}$\ldots

\uline{Remark} We have not for the moment taken account of the congruence
conditions
\[
\begin{pmatrix}a&b\\c&d\end{pmatrix}\equiv1\ (2)
\qquad
\begin{pmatrix}r&s\\t&u\end{pmatrix}\equiv1\ (\unsure{2})
\]
which express that these matrices arise modulo signs by ``reduction''
from elements of
$\unsure{\widehat{\widetilde{\Pi}}_{0,3}}$, $\widehat{\Pi}_{0,3}$.
For the quantities $(c,d)$ \deleted{\unsure{fixed}} one finds the
consequences
\[
\left\{\begin{aligned}
d&\equiv\pm1\ (4)\quad\text{hence }d^2\equiv1\ (4)\\
c&\equiv0,2\ (4)\quad\text{hence }c^2\equiv0\ (4)
\end{aligned}\right.
\ \Longrightarrow\
\begin{aligned}
c^2+d^2+cd&\equiv1+cd\ (4)\\
&\equiv1+c\ (4)
\end{aligned}
\]
hence
\[
c^2+d^2+cd\equiv1+c\ \deleted{\unsure{(4)}}(4)
\quad\deleted{\unsure{\Longleftrightarrow}\ \unsure{cd}\ cd\equiv0\ (4)}
\quad\text{hence}
\]
hence
\[
\mu=\varepsilon_3(c^2+cd+d^2)\equiv
\deleted{\unsure{\Longleftrightarrow c\equiv0\ (4)}}
\varepsilon_3(1+c)\ (4)
\]
\deleted{(N.B. $d$ is invertible modulo $4$ since $d\equiv1\ (2)$),}

\margin{165}

\scanpage{167}
\mspage{382}

We thus see that knowledge of $c$ \deleted{\illegible{1}}
\inserted{modulo $4$}, and that of the sign $ad-bc=\varepsilon_3$,
allows one to find the sign $\varepsilon_2$ \deleted{\unsure{by}} by
\[
\deleted{\unsure{\varepsilon_2}\equiv\varepsilon_3\unsure{c}\quad
\varepsilon_2\equiv\varepsilon_3(1+c)\ (4)}
\]
\[
\tag{90}
\varepsilon_2\equiv\varepsilon_3(1+c)\ (4)
\quad\text{i.e.}\quad
\varepsilon_2\varepsilon_3\equiv1+c\ (4)
\]
If one sets out, for $c,d,t,u$ fixed, satisfying
\[
\tag{91}
\left\{\begin{aligned}
c&\equiv t\equiv0\ (2)\qquad d\equiv u\equiv1\ (2)\\
(c^2+\unsure{cd}+d^2)&\unsure{\equiv|=}\deleted{\unsure{(1+c}}
\varepsilon(t^2+u^2)\in\widehat{\mathbb Z}^{*}
\end{aligned}\right.
\]
where $\varepsilon$ \deleted{\illegible{1}}\inserted{$\in\{\pm1\}$}
is the sign defined by
\[
\varepsilon\equiv1+c\ (4)
\]
($\varepsilon=1$ if $c\equiv0\ (4)$, $\varepsilon=-1$ if
$c\equiv2\ (4)$), to find the corresponding $a,b,r,s$,
satisfying conditions (89), and moreover the congruence conditions
\[
\tag{92}
\unsure{b}\equiv\unsure{s}\equiv0\ (2)\qquad
a\equiv r\equiv1\ (2),
\]
one is led to lift an already given solution modulo $2$ of this system
of equations to a solution. In terms of the \inserted{$\widehat{\mathbb Z}$-}
torsor $\unsure{\mathcal T}$ of solutions, this means that one gives oneself
an element of the corresponding quotient $\mathbb Z/2\mathbb Z$-torsor
--- the set of its liftings is identified with a torsor under
$2\widehat{\mathbb Z}\simeq\widehat{\mathbb Z}$\ldots{} No mystery!

\margin{166}

\scanpage{168}

\uline{Remark} The additive group $\widehat{\mathbb Z}$ operates on the set
of solutions
$(\alpha_0,f)=\left(\begin{pmatrix}a&b\\c&d\end{pmatrix},
\begin{pmatrix}r&s\\t&u\end{pmatrix}\right)$
of the equations under consideration (48)(49)(56)
\[
\tag{93}
\det\alpha_0\in\{\pm1\},\quad\det f=1,\quad
\alpha_0\rho(\alpha_0)^{-1}
=\varepsilon\ell_0^{\nu}\bigl(f\sigma_\infty(f)^{-1}\bigr),
\quad\varepsilon\in\{\pm1\},
\]
by
\[
\tag{\unsure{93|94}}
\begin{aligned}
(\alpha_0,f)&\longmapsto(\alpha'_0,f')
=(\ell_0^n\alpha_0,\ell_0^nf)\quad\deleted{=(\unsure{a}-2nc,\ b-}\\
&=\left(\begin{pmatrix}a-2nc&b-2nd\\c&d\end{pmatrix},
\begin{pmatrix}r-2nt&s-2nu\\t&u\end{pmatrix}\right)
\end{aligned}
\]
($n\in\widehat{\mathbb Z}$).
Up to the factor $2$ in $2n$, this is precisely the indeterminacy in
\deleted{\illegible{1}} the set of solutions \inserted{of (93)}, when
$c,d,t,u$ are fixed (cf. (87) and the proposition that goes with it).
\inserted{(the factor $2$ is \unsure{ensured} \unsure{only when}
\unsure{one considers|in considering} solutions satisfying
$\alpha'_0\equiv f'_0\equiv\mathrm{Id}\ (2)$)}
One can \deleted{\unsure{thus}} say that what is determined
intrinsically \inserted{in
$\unsure{\mathrm{Gl}(2,\widehat{\mathbb Z})^{\pm1\prime}|\mathrm{Gl}(2,\widehat{\mathbb Z})^{\pm1}}$}
in terms of the \inserted{loop} automorphism $\uline{u}$ of
$\widehat{\Pi}_{0,3}$ commuting with $\sigma_\infty$, is neither more nor less
than \deleted{the \unsure{values}}
\[
\tag{95}
\begin{aligned}
\text{a)}\quad&\text{the multiplier }\mu\in\widehat{\mathbb Z}^{*}
\end{aligned}
\]
\deleted{\unsure{which gives} $\varepsilon_2=\unsure{\varepsilon_2|\xi_2}(\mu),\
\varepsilon_3=\unsure{\varepsilon_3|\xi_3}(\mu)$}
\inserted{(\unsure{hence|whence} $\unsure{\varepsilon_2}=\unsure{\varepsilon_2|\xi_2}(\mu),\
\varepsilon_3=\unsure{\varepsilon_3|\xi_3}(\mu)$)}

b) the quantities
\deleted{$\unsure{\delta}=\unsure{\varepsilon_3},\mu,\
c,d,t,u\in\widehat{\mathbb Z}$, \unsure{satisfying}}
the $c,d,t,u\in\widehat{\mathbb Z}$ satisfying the
everlasting relation (86)
\[
\tag{96}
\varepsilon_2(t^2+u^2)=\varepsilon_2(c^2+cd+d^2)=1/\mu
\]
and the congruence relations
\[
\tag{97}
c\equiv t\equiv0\ (2)\qquad d\equiv u\equiv1\ (2)
\]

It is true that we have seen that there is a \inserted{natural} way in
$\widehat{\Pi}_{0,3}$ to normalize the choice of the $f$, so as to remove
the \unsure{ambiguity|ambiguity} of the $\ell_0^nf$
(cf. \S\,41 \inserted{(24bis)} where this was touched upon --- I will come back to it
again later\ldots) but it seems to me \inserted{\uline{not very}}
interesting

\margin{167}